% Revize výukové přípravy, výchozí soubor GitHub Work-MAT, 86. hodina.
% Zdroje: původní hodina 86; Realisticky.cz 020203–020207, 020310–020311
% Úlohy jsou vlastní / úlohy výchozí přípravy, případně nově didakticky obměněné.
\documentclass[11pt,a4paper]{article}
\input{../../sablony/pripravy/priprava.sty}
\usepackage{tikz}
\nastavPripravu{1. ročník}{Rovnice a nerovnice}{86}{Průběžný test I}
\begin{document}
\blok{Průběžný test I – lineární rovnice, nerovnice a soustavy}
Test probíhá po dokončení hodin 72--85. Žáci obdrží \textbf{samostatnou variantu A nebo B}; tuto učitelskou přípravu s řešeními nedostávají.
\blok{Organizace a hodnocení}
\begin{itemize}
\item Pokyny 3 min, samostatná práce 22 min, navazující individuální oprava či procvičování 15 min, odevzdání a závěr 5 min.
\item Čtyři úlohy po 3, 3, 4 a 3 bodech; celkem \textbf{13 bodů}. Hodnoť i nezáporné / záporné násobitele, intervaly a zkoušku v soustavách.
\item Zadání A/B a samostatné učitelské klíče jsou ve složce \texttt{1\_rocnik/pracovni\_listy/Rovnice/}.
\end{itemize}
\blok{Ukázka rozboru po testu}
\begin{prikladReseniBox}[title={Příklad 1}]
Zkontrolujte, zda se řešení $-2(x-1)\ge6$ správně píše $x\ge-2$.
\tcblower
\[\begin{aligned}
-2x+2&\ge6 \\
-2x&\ge4 \\
x&\le-2 \\
K&=(-\infty;-2\rangle
\end{aligned}\]
Znaménko se při dělení číslem $-2$ obrací.
\end{prikladReseniBox}
\begin{zadaniReseniBox}[title={Samostatná práce 1}]
Vyřešte soustavu a výsledek zapište jako množinu uspořádaných dvojic: $\begin{cases}2x+y=9,\\x-y=0.\end{cases}$
\tcblower
\[\begin{aligned}
y&=x \\
2x+x&=9 \\
x&=3,\quad y=3 \\
K&=\{[3;3]\}
\end{aligned}\]
Zkouška: $2\cdot3+3=9$, $3-3=0$.
\end{zadaniReseniBox}
\poznamka{Sleduj typické chyby anonymně: dělení záporným číslem, průnik intervalů, chybná eliminace a neuspořádaný zápis řešení soustavy.}
\blok{Domácí úkol}
\begin{zadaniReseniBox}[title={Domácí úkol}]
\begin{enumerate}[label=\alph*),itemsep=2pt]
\item Vysvětlete chybu v úpravě $-4x<12\Rightarrow x<-3$.
\item $3(x-2)=2x+1$.
\item $\dfrac{2x-1}{3}\ge\dfrac{x+2}{2}$.
\item $\begin{cases}x+y=6,\\x-y=4.\end{cases}$
\item $x+1\le4$ a zároveň $x\ge0$.
\end{enumerate}
\tcblower
\begin{enumerate}[label=\alph*),itemsep=1pt]
\item
\[\begin{aligned}
-4x&<12 \\
x&>-3 \\
K&=(-3;+\infty)
\end{aligned}\]\par Po vydělení záporným číslem se obrátí znak nerovnosti.
\item
\[\begin{aligned}
3x-6&=2x+1 \\
x&=7 \\
K&=\{7\}
\end{aligned}\]
\item
\[\begin{aligned}
2(2x-1)&\ge3(x+2) \\
x&\ge8 \\
K&=\langle8;+\infty)
\end{aligned}\]
\item
\[\begin{aligned}
2x&=10 \\
x&=5,\ y=1 \\
K&=\{[5;1]\}
\end{aligned}\]
\item
\[\begin{aligned}
x&\le3 \\
x&\ge0 \\
K&=\langle0;3\rangle
\end{aligned}\]
\end{enumerate}
\end{zadaniReseniBox}
\poznamka{Po testu ověř, zda žák rozumí svému postupu, nikoli jen výsledku. Doporučená kontrola: u dvou různých řešení na tabuli žák identifikuje první chybný řádek.}
\end{document}
